% language=us

\environment context-2025-style


\setuptyping
  [style=\switchtobodyfont[7pt,tt]]

\definelayer
  [extratext]
  [preset=rightbottom,
   height=\textheight,
   width=\textwidth]

\setupbackgrounds
  [text]
  [background=extratext]

% \definehighlight[notabene][style=bf]
\definehighlight[notabene][style=\glyphweight100\underbar]

\setuptyping
  [align=hangright,
   keeptogether=paragraph,
   numbering=]

\startdocument
  [title={Signals},
%    banner={keeping track},
   banner={Willi, Mojca & Hans},
   location={context\enspace{\bf 2025}\enspace meeting}]

\starttitle[title=From beamer to lamps]

\startitemize
\startitem
    We wanted Tove's Bacho\TEX\ \quote {flags} presentation to have some weird
    properties.
\stopitem
\startitem
    I was working on a large 65 inch screen on 3 meter distance due to an eye
    issue.
\stopitem
\startitem
    I had some hue color lamps laying around unused and the home|/|office lighting
    is driven by a few hue hubs already.
\stopitem
\startitem
    Control of lamps is actually done by a \LUA\ script (server) run by \LUATEX,
    a way to overcome limitations (and performance issues).
\stopitem
\startitem
    So, controlling lamps was relatively easy. Making them respond to a change in
    \PDF\ page not that hard.
\stopitem
\startitem
    That in return triggered using lamps to indicate a \CONTEXT\ run state and
    result.
\stopitem
\stopitemize

\stoptitle

\starttitle[title=From lamps to leds]

\startitemize
\startitem
    At Bacho\TEX\ we played a bit with a led strip (on a hue led controller)
    that Mojca gave me.
\stopitem
\startitem
    That demonstrated that instead of (large) lamps we could go smaller.
\stopitem
\startitem
    We used 30 leds and therefore could play with configurations. Eventually
    24 leds were the optimum (1 2 3 4 6 8 12).
\stopitem
\startitem
    By putting them in a circle one gets a kind of clock (watch).
\stopitem
\startitem
    We use colors: blue (busy), yellow (current), red (error), orange (problem),
    green (okay).
\stopitem
\startitem
    We can work on segments (24/8), quadrants (24/4), or squids (24/1). The first
    two report per run (if more are needed), the later the current run with page
    indicator.
\stopitem
\stopitemize

\stoptitle

\starttitle[title=From leds to gadget]

\startitemize
\startitem
    Mojca suggested to make a dedicated device, and gave me a Raspberry Pico
    to play with.
\stopitem
\startitem
    The quadrant, segment, squid model was \quote {perfected} in an experimental
    setup.
\stopitem
\startitem
    We still support hue, we can also respond to http requests but the focus is
    now on fast serial communication.
\stopitem
\startitem
    Controlling the gadget (aka \CONTEXT\ watch, as it is round) happens from
    the \type {context} script and in squid mode also from the shipout routine.
\stopitem
\startitem
    Serial communication is built into \LUAMETATEX, and it works okay on
    \WINDOWS, \LINUX\ and \OSX.
\stopitem
\stopitemize

\stoptitle

\starttitle[title=From gadget to tool]

\startitemize
\startitem
    We can see progress and result state using quadrant, segment, or squid mode.
    The documents top line indicates what has to be used.
\stopitem
\startitem
    In squid mode we can also get feedback about e.g. overfull boxes, missing
    characters, options that slow down a run, or potentially interfering
    third party code.
\stopitem
\startitem
    In segment mode we can show the state of 8 parallel jobs, this is used when
    running the test suite.
\stopitem
\startitem
    You can control the device by sending commands and thereby integrate it
    in workflows. We use this when generating the distribution.
\stopitem
\startitem
    Future versions can have a clock mode. We already can count (numbers) and do
    some simple text.
\stopitem
\stopitemize

\stoptitle

\starttitle[title=How it fits in]

\startitemize
\startitem
    It all fits in a feedback|/|runner model. When a run goes bad, one gets
    an error document.
\stopitem
\startitem
    But instead one can get a QR code as result, encoding the error (there
    is a limited set so basically they are pictograms.
\stopitem
\startitem
    Using leds (color and/or position) in practice is a faster feedback than
    watching and reading a log file.
\stopitem
\startitem
    It also marks the 20\high{th} anniversary of \LUATEX, we've come a long
    way and still finds ways to extend it.
\stopitem
\startitem
    The software is there but can be upgraded as we go. There's also a manual
    and a helper script in the distribution.
\stopitem
\startitem
    As the gadget can be controlled by a script you can also use it otherwise,
    but as usual, we welcome ideas by users.
\stopitem
\stopitemize

\stoptitle

\starttitle[title=Dejavu]

\startitemize
\startitem
I played a lot with a RCA 1802 dev kit when was a kid myself. It was one of nicer
cpu designs, quit symmetrical with static memory and single step capabilities.
Think of the kind of chips that currently travel the edges of our solar system.
\stopitem
\startitem
Writing games and such in 512 bytes mem (or was it 256) is a challenge.
\stopitem
\startitem
Inventing wheels is part of it. One learns a lot: interrupts, self adapting
functions, figuring out how to save to and read from tape, dma and sensors.
\stopitem
\startitem
When we started with \LUATEX, at some point Mojca send me a esp micro controller
running \LUA. It brought back memories.
\stopitem
\startitem
Years later Harald made a laser cut \CONTEXT\ logo controlled by a rpi and later
Willi made an enclosure with a next generation embedded esp micro controller.
\stopitem
\startitem
Today we show a gadget that has yet another generation of that esp (although we
actually started with a rpi pico).
\stopitem

\blank[2*big]

\startitem
For a moment think back to when \TEX\ was written: this tiny chip has more cpu
cycles, flash, plenty static and optionally extra memory and quite some IO
features, it can run original \TEX\ faster than the machines on which
he program was written.
\stopitem
\stopitemize

\stoptitle

\starttitle[title=The workshop]

\startplacefigure[location={left,none}]
\externalfigure[release-willi-watch.jpg][height=.7th]
\stopplacefigure

\startitemize
\startitem
    If there is interest we will make more boards that can be mounted in a clock
    encasing with a 3D printed enclosure.
\stopitem
\startitem
    But \unknown\ at this meeting you will get the collectors item: a high
    quality wooden Willi designed stand with the first official hard- and
    software release.
\stopitem
\startitem
    But you need to work a bit on it, welcome to:

    {\em The traditional \CONTEXT\ workshop event!}
\stopitem
\stopitemize

\stoptitle

\stopdocument


% context/signal/testthem : configure hue 3 or 4
